% !TEX program = xelatex
\documentclass[14pt]{extarticle}

% ============================================================
% GEOMETRY
% ============================================================
\usepackage[
  a4paper,
  left=2.5cm,
  right=1.5cm,
  top=2cm,
  bottom=2cm
]{geometry}

% ============================================================
% МОВА ТА ШРИФТИ
% ============================================================
\usepackage[ukrainian]{babel}
\usepackage{fontspec}
\setmainfont{Times New Roman}

% ============================================================
% ОСНОВНИЙ ТЕКСТ
% ============================================================
\usepackage{setspace}
\onehalfspacing
\setlength{\parindent}{1.27cm}
\setlength{\parskip}{0pt}
\usepackage{microtype}

% ============================================================
% ЗАГОЛОВКИ
% ============================================================
\usepackage{titlesec}

\titleformat{\section}
  {\bfseries\fontsize{16pt}{24pt}\selectfont\centering}
  {\thesection}{0.5em}{}
\titlespacing*{\section}{0pt}{20pt}{6pt}

\titleformat{\subsection}
  {\bfseries\fontsize{14pt}{21pt}\selectfont}
  {\thesubsection}{0.5em}{}
\titlespacing*{\subsection}{1.27cm}{6pt}{6pt}

% ============================================================
% СПИСКИ
% ============================================================
\usepackage{enumitem}

\setlist[itemize,1]{
  label=$\bullet$,
  leftmargin=1.27cm,
  labelwidth=0.5cm,
  labelsep=0.3cm,
  itemindent=0pt,
  itemsep=3pt,
  topsep=6pt,
  parsep=0pt,
  partopsep=0pt
}

\setlist[enumerate,1]{
  label=\arabic*.,
  leftmargin=2.54cm,
  labelwidth=0.63cm,
  labelsep=0.54cm,
  itemindent=0pt,
  itemsep=0pt,
  topsep=0pt,
  parsep=0pt,
  partopsep=0pt
}

% ============================================================
% ТАБЛИЦІ
% ============================================================
\usepackage{array}
\usepackage{longtable}
\usepackage{booktabs}
\usepackage{tabularx}

% ============================================================
% ГІПЕРПОСИЛАННЯ
% ============================================================
\usepackage[hidelinks,breaklinks=true]{hyperref}

\begin{document}

\sloppy
\emergencystretch=3em

% ==================== ТИТУЛЬНА ЧАСТИНА ====================
\begin{center}
\textbf{Лабораторна робота №2}

\medskip

\textbf{\MakeUppercase{Особливості побудови процесу управління проєктами (формування концептуальних понять проєкту). Обґрунтування сфери застосування та об'єму робіт проєкту.}}
\end{center}

\medskip

\textbf{Мета роботи:} ознайомитись з основними процесами управління проєктами, особливостями побудови процесу управління проєктом та роллю менеджера в цих процесах. Навчитись окреслювати основні та додаткові етапи роботи.

\textbf{Проєкт:} <<Система для автоматичної нормалізації суржику та діалектних форм української мови>>.

% ==================== РОЗДІЛ 1: КОНЦЕПТУАЛЬНІ ПОНЯТТЯ ====================
\section*{1. Формування концептуальних понять проєкту}
\addcontentsline{toc}{section}{1. Формування концептуальних понять проєкту}

\subsection*{1.1. Визначення проєкту}

Проєкт <<Система для автоматичної нормалізації суржику та діалектних форм української мови>> - це тимчасове зусилля, спрямоване на створення унікального програмного продукту: нейромережевої системи, що перетворює тексти, написані українськими діалектами (гуцульським, бойківським, закарпатським) та російсько-українським суржиком, на літературну українську мову. Проєкт реалізується в рамках бакалаврської кваліфікаційної роботи за спеціальністю 122 <<Комп'ютерні науки>> у Національному університеті <<Львівська політехніка>>.

\subsection*{1.2. Команда проєкту та ролі}

\textbf{Керівник проєкту (виконавець):} студент групи ШІ-42 Коваль Денис Петрович. Відповідає за планування, виконання всіх технічних етапів, контроль якості та оформлення результатів. Оскільки проєкт є індивідуальною бакалаврською роботою, виконавець суміщає ролі менеджера проєкту, розробника та аналітика даних.

\textbf{Науковий керівник (спонсор):} Сивоконь Олексій Олегович, асистент кафедри систем штучного інтелекту, ІКНІ. Забезпечує методичне спрямування, контроль відповідності дослідження академічним стандартам, рецензування результатів. У термінології PMBOK виконує роль спонсора, який <<розчищає шлях>> у питаннях доступу до ресурсів та консультацій.

\subsection*{1.3. Ромб конкурентних обмежень}

Ромб конкурентних обмежень проєкту визначається чотирма вершинами:

\textbf{Кошти.} Проєкт є академічним і не має виділеного фінансового бюджету. Основні витрати - обчислювальні ресурси: baseline тренується на CPU (безкоштовно), донавчання LLM потребує GPU з $\geq$24~GB відеопам'яті (QLoRA знижує вимоги). Використовуються відкриті дані та моделі.

\textbf{Час.} Обмежений одним навчальним семестром з фіксованою датою захисту бакалаврської роботи.

\textbf{Якість.} Вимірюється метриками BLEU, chrF++, TER, COMET на контрольних наборах з вручну анотованих пар. Baseline для гуцульського діалекту досягає BLEU~$\approx$51.

\textbf{Результат.} Працююча система нормалізації з веб-інтерфейсом, що підтримує 3 діалекти та суржик, з документованими результатами експериментів.

% ==================== РОЗДІЛ 2: ЕТАПИ ПРОЄКТУ ====================
\section*{2. Етапи проєкту та контрольні списки (checklists)}
\addcontentsline{toc}{section}{2. Етапи проєкту та контрольні списки}

Відповідно до процесу управління проєктами (PMBOK), проєкт проходить п'ять етапів: ініціювання, планування, виконання, контроль та закриття. Для кожного етапу сформовано опорний технічний контрольний список (checklist) із очікуваними результатами.

\subsection*{2.1. Ініціювання}

На етапі ініціювання визначається доцільність проєкту та формується його загальне бачення. Основні результати:

\begin{longtable}{|p{0.55\textwidth}|p{0.35\textwidth}|}
\hline
\textbf{Checklist-елемент} & \textbf{Результат} \\
\hline
\endfirsthead
\hline
\textbf{Checklist-елемент} & \textbf{Результат} \\
\hline
\endhead
Визначення проблеми (відсутність інструменту нормалізації діалектів/суржику) & Опис проблеми та актуальності \\
\hline
Визначення сфери застосування (7 базових питань) & Документ сфери застосування \\
\hline
Формулювання мети за SMART & Конкретна, вимірювана мета \\
\hline
Визначення зацікавлених сторін & Перелік стейкхолдерів \\
\hline
Узгодження теми з науковим керівником & Затверджена тема БКР \\
\hline
Попередній огляд доступних ресурсів (дані, моделі) & Перелік ключових джерел \\
\hline
\end{longtable}

\subsection*{2.2. Планування}

На етапі планування деталізуються роботи, визначаються ресурси, ризики та розклад. Основні результати:

\begin{longtable}{|p{0.55\textwidth}|p{0.35\textwidth}|}
\hline
\textbf{Checklist-елемент} & \textbf{Результат} \\
\hline
\endfirsthead
\hline
\textbf{Checklist-елемент} & \textbf{Результат} \\
\hline
\endhead
Огляд літератури (NMT, GEC, діалектна нормалізація, суржик) & Аналітичний огляд джерел \\
\hline
Аналіз наборів даних (ручний корпус, синтетика, словник, Spivavtor) & Опис джерел даних та їх ролей \\
\hline
Вибір архітектур (seq2seq baseline, Lapa LLM з QLoRA) & Обґрунтований вибір моделей \\
\hline
Визначення метрик якості (BLEU, chrF++, TER, COMET) & Набір метрик з цільовими значеннями \\
\hline
Визначення технологічного стеку (Python, PyTorch, HF Transformers, PEFT, Gradio) & Перелік інструментів \\
\hline
Формування розкладу етапів & План-графік робіт \\
\hline
Ідентифікація ризиків & Реєстр ризиків \\
\hline
Визначення вимог до обчислювальних ресурсів & Специфікація (CPU для baseline, GPU $\geq$24~GB для LLM) \\
\hline
\end{longtable}

\subsection*{2.3. Виконання}

Етап виконання - це безпосередня реалізація запланованих робіт. Основні результати:

\begin{longtable}{|p{0.55\textwidth}|p{0.35\textwidth}|}
\hline
\textbf{Checklist-елемент} & \textbf{Результат} \\
\hline
\endfirsthead
\hline
\textbf{Checklist-елемент} & \textbf{Результат} \\
\hline
\endhead
Збір та підготовка даних для гуцульського діалекту & Паралельний корпус (ручний + синтетичний), словник \\
\hline
Побудова та навчання baseline seq2seq моделі & Навчена модель, BLEU~$\approx$51 \\
\hline
Аналіз помилок baseline & Звіт про типові помилки \\
\hline
Формування мультидіалектного корпусу (бойківський, закарпатський, суржик) & Розширений корпус із синтетикою на базі Spivavtor \\
\hline
Синтетична генерація пар (словниково-керована, правилова, LLM з RAG) & Набори синтетичних пар із контролем якості \\
\hline
Інструкційне донавчання Lapa LLM з QLoRA & Донавчена модель \\
\hline
Оцінювання якості на контрольних наборах & Таблиця метрик по діалектах/суржику \\
\hline
Розробка веб-інтерфейсу (Gradio) & Працюючий демо-інтерфейс \\
\hline
\end{longtable}

\subsection*{2.4. Контроль}

Етап контролю відбувається паралельно з виконанням і забезпечує відстеження прогресу та реагування на відхилення. Основні результати:

\begin{longtable}{|p{0.55\textwidth}|p{0.35\textwidth}|}
\hline
\textbf{Checklist-елемент} & \textbf{Результат} \\
\hline
\endfirsthead
\hline
\textbf{Checklist-елемент} & \textbf{Результат} \\
\hline
\endhead
Відстеження відповідності розкладу & Звіт про дотримання термінів \\
\hline
Моніторинг метрик якості на validation-наборах під час навчання & Графіки loss, BLEU по епохах \\
\hline
Перевірка якості синтетичних даних (фільтрація, аудит) & Лог відфільтрованих прикладів \\
\hline
Контроль розбиття train/validation/test (ручні дані лише у val/test) & Підтвердження коректності розбиття \\
\hline
Відстеження змін у вимогах або підході & Журнал змін \\
\hline
Регулярні консультації з науковим керівником & Протоколи обговорень \\
\hline
\end{longtable}

\subsection*{2.5. Закриття}

На етапі закриття фіксуються результати, оформлюється документація та передається досвід. Основні результати:

\begin{longtable}{|p{0.55\textwidth}|p{0.35\textwidth}|}
\hline
\textbf{Checklist-елемент} & \textbf{Результат} \\
\hline
\endfirsthead
\hline
\textbf{Checklist-елемент} & \textbf{Результат} \\
\hline
\endhead
Оформлення тексту бакалаврської кваліфікаційної роботи & Готовий документ БКР \\
\hline
Фіксація фінальних метрик якості по всіх діалектах та суржику & Таблиця результатів \\
\hline
Публікація коду та даних у репозиторії & Відкритий репозиторій \\
\hline
Підготовка презентації до захисту & Слайди \\
\hline
Фіксація отриманого досвіду (lessons learned) & Розділ висновків у БКР \\
\hline
Захист бакалаврської роботи & Оцінка \\
\hline
\end{longtable}

% ==================== РОЗДІЛ 3: РИЗИКИ ====================
\section*{3. Ризики та перешкоди}
\addcontentsline{toc}{section}{3. Ризики та перешкоди}

Нижче подано ключові ризики, які можуть виникнути при розробленні проєкту, та заходи з їх мінімізації.

\begin{longtable}{|p{0.04\textwidth}|p{0.28\textwidth}|p{0.18\textwidth}|p{0.36\textwidth}|}
\hline
\textbf{№} & \textbf{Ризик} & \textbf{Ймовірність / Вплив} & \textbf{Заходи мінімізації} \\
\hline
\endfirsthead
\hline
\textbf{№} & \textbf{Ризик} & \textbf{Ймовірність / Вплив} & \textbf{Заходи мінімізації} \\
\hline
\endhead
1 & Недостатній обсяг ручних даних для діалектів (крім гуцульського) & Висока / Високий & Синтетичне розширення на базі Spivavtor, словникових правил та LLM-генерації з RAG-підказками \\
\hline
2 & Низька якість синтетичних даних (нереалістичні конструкції) & Середня / Високий & Фільтрація дублікатів, контроль довжин, словникова валідація, вибірковий ручний аудит \\
\hline
3 & Недостатні обчислювальні ресурси для донавчання LLM (12B параметрів) & Середня / Високий & Використання QLoRA для зменшення вимог до відеопам'яті; fallback на менші моделі \\
\hline
4 & Затримка на етапі підготовки даних, що скорочує час на донавчання & Середня / Середній & Паралельне формування корпусів для різних діалектів; чіткий план-графік з контрольними точками \\
\hline
5 & Catastrophic forgetting при мультидіалектному навчанні & Середня / Середній & Стратифіковане балансування семплінгу за діалектами, окремі контрольні набори для кожного різновиду \\
\hline
6 & Метрики не відображають реальну якість нормалізації & Низька / Середній & Використання нейронних метрик (COMET) поряд із класичними; вибіркова ручна оцінка \\
\hline
7 & Зміна або недоступність зовнішніх ресурсів (моделі, датасети на HuggingFace) & Низька / Високий & Локальне кешування всіх ключових даних та моделей на початку проєкту \\
\hline
8 & Невідповідність результатів академічним вимогам до оформлення & Низька / Середній & Регулярні консультації з науковим керівником, дотримання шаблону оформлення \\
\hline
\end{longtable}

% ==================== ВИСНОВКИ ====================
\section*{Висновки}
\addcontentsline{toc}{section}{Висновки}

У ході виконання лабораторної роботи було сформовано концептуальні поняття проєкту: визначено його як тимчасове зусилля зі створення унікального програмного продукту, описано команду (виконавець та науковий керівник) та ромб конкурентних обмежень (кошти, час, якість, результат).

Для кожного з п'яти етапів управління проєктом (ініціювання, планування, виконання, контроль, закриття) розроблено технічні контрольні списки (checklists) з конкретними елементами та очікуваними результатами. Контрольні списки прив'язані до реальних артефактів проєкту: від корпусів даних і навчених моделей до метрик якості та демо-інтерфейсу.

Ідентифіковано 8 ключових ризиків проєкту з оцінкою ймовірності та впливу, для кожного запропоновано конкретні заходи мінімізації. Найбільш критичними є: дефіцит ручних даних для діалектів (мінімізується синтетичним розширенням), обмеженість обчислювальних ресурсів (мінімізується QLoRA) та ризик затримки на підготовці даних (мінімізується паралельною роботою та контрольними точками).

\end{document}
